\documentclass[11pt,a4paper]{article}
\usepackage[utf8]{inputenc}
\usepackage[T1]{fontenc}
\usepackage[french]{babel}
\usepackage{amsmath,amssymb,amsfonts,amsthm}
\usepackage{geometry}
\geometry{hmargin=2.5cm,vmargin=3cm}
\usepackage{enumitem}
\usepackage{tabularx}
\usepackage{booktabs}

\title{\textbf{Corrigé du Sujet d'Entraînement n°1 : Statistiques à une variable}\\\large Préparation au CAPES de Mathématiques}
\author{Correction détaillée et commentaires pédagogiques}
\date{Session 2025-2026}

\begin{document}

\maketitle

\section*{Exercice 1 : Questions de cours (Démonstrations fondatrices)}

\begin{enumerate}
    \item \textbf{Propriété de moindre carré de la moyenne}
    \begin{enumerate}
        \item Soit $\theta \in \mathbb{R}$. Développons le terme général sous la somme de la fonction $g$ :
        \[ g(\theta) = \sum_{i=1}^{k} n_{i}(x_{i}-\theta)^{2} = \sum_{i=1}^{k} n_{i}(x_{i}^{2} - 2\theta x_{i} + \theta^{2}) \]
        Par linéarité de la sommation, on sépare l'expression en trois sommes distinctes et on extrait les facteurs constants par rapport à l'indice $i$ :
        \[ g(\theta) = \sum_{i=1}^{k} n_{i}x_{i}^{2} - 2\theta \sum_{i=1}^{k} n_{i}x_{i} + \theta^{2} \sum_{i=1}^{k} n_{i} \]
        En utilisant les relations fondamentales du cours :
        \begin{itemize}
            \item $\sum_{i=1}^{k} n_{i} = N$ (effectif total)
            \item $\sum_{i=1}^{k} n_{i}x_{i} = N\overline{x}$ (par définition de la moyenne $\overline{x} = \frac{1}{N}\sum_{i=1}^{k}n_ix_i$)
        \end{itemize}
        On obtient l'expression polynomiale en $\theta$ recherchée :
        \[ g(\theta) = N\theta^{2} - 2N\overline{x}\theta + \sum_{i=1}^{k} n_{i}x_{i}^{2} \]
        C'est un polynôme du second degré de la forme $A\theta^2 + B\theta + C$ avec :
        \[ A = N \quad , \quad B = -2N\overline{x} \quad , \quad C = \sum_{i=1}^{k} n_{i}x_{i}^{2} \]
        
        \item Le coefficient $A = N$ est strictement positif ($N \in \mathbb{N}^*$). La fonction polynomiale $g$ est donc strictement convexe et admet un unique minimum global sur $\mathbb{R}$.
        Ce minimum est atteint au sommet de la parabole, soit en :
        \[ \theta_0 = -\frac{B}{2A} = -\frac{-2N\overline{x}}{2N} = \overline{x} \]
        Ainsi, la fonction $g$ atteint son minimum global unique en $\theta = \overline{x}$.
        
        \item \textbf{Interprétation géométrique/intuitive :} Le terme $(x_i - \theta)^2$ représente le carré de la distance euclidienne entre une valeur de la série et un paramètre central $\theta$. Minimiser la somme de ces carrés revient à chercher le point $\theta$ « le plus proche » de toutes les données au sens de la distance quadratique. La moyenne est donc le centre de gravité (ou barycentre) de la distribution.
    \end{enumerate}

    \item \textbf{Optimisation en valeur absolue et médiane}
    \begin{enumerate}
        \item La fonction $h(\theta) = \sum_{i=1}^{N} |x_{i}-\theta|$ est une somme de fonctions de la forme $\theta \mapsto |x_i - \theta|$. Chacune de ces fonctions est continue sur $\mathbb{R}$ et affine par morceaux (linéaire de pente $-1$ si $\theta < x_i$ et de pente $+1$ si $\theta > x_i$). Par somme de fonctions ayant ces propriétés, $h$ est continue, affine par morceaux et convexe sur $\mathbb{R}$.
        
        \item Soit $\theta \notin \{x_1, \dots, x_N\}$. La fonction valeur absolue est dérivable partout sauf en $0$. Donc $h$ est dérivable en $\theta$. Pour chaque $i \in \{1, \dots, N\}$ :
        \begin{itemize}
            \item Si $x_i < \theta$, alors $|x_i - \theta| = \theta - x_i$, de dérivée par rapport à $\theta$ égale à $1$.
            \item Si $x_i > \theta$, alors $|x_i - \theta| = x_i - \theta$, de dérivée par rapport à $\theta$ égale à $-1$.
        \end{itemize}
        On en déduit que la dérivée $h'(\theta)$ vaut :
        \[ h'(\theta) = \sum_{i=1}^{N} \text{sgn}(\theta - x_i) = \text{card}(\{i \ / \ x_i < \theta\}) - \text{card}(\{i \ / \ x_i > \theta\}) \]
        Autrement dit, la dérivée en un point correspond au nombre de données situées à gauche de $\theta$ moins le nombre de données situées à droite de $\theta$.
        
        \item Posons $N = 2p+1$. La série étant ordonnée, la médiane unique au sens pratique est $M_e = x_{p+1}$. Étudions le signe de $h'(\theta)$ sur les intervalles ouverts délimités par les valeurs distinctes :
        \begin{itemize}
            \item \textbf{Si $\theta < M_e = x_{p+1}$} : $\theta$ est situé avant la valeur d'indice $p+1$. Le nombre de données inférieures à $\theta$ est au maximum $p$, tandis que le nombre de données supérieures à $\theta$ est au moins $p+1$. Dès lors, $\text{card}(\{i \ / \ x_i < \theta\}) < \text{card}(\{i \ / \ x_i > \theta\})$, d'où $h'(\theta) < 0$. La fonction $h$ est strictement décroissante sur $]-\infty ; M_e[$.
            \item \textbf{Si $\theta > M_e = x_{p+1}$} : Par un raisonnement symétrique, le nombre de données situées à gauche de $\theta$ est au moins $p+1$ et à droite au plus $p$. On a $\text{card}(\{i \ / \ x_i < \theta\}) > \text{card}(\{i \ / \ x_i > \theta\})$, d'où $h'(\theta) > 0$. La fonction $h$ est strictement croissante sur $]M_e ; +\infty[$.
        \end{itemize}
        Par continuité de $h$, la fonction décroît strictement jusqu'en $M_e$ puis croît strictement après $M_e$. Le minimum global unique de $h$ est donc bien atteint en la médiane $\theta = M_e$.
    \end{enumerate}
\end{enumerate}

\newpage

\section*{Exercice 2 : Effets de transformations affines et indicateurs robustes}

\begin{enumerate}
    \item \textbf{Harmonisation par translation et dilatation ($Y = 1,2X + 1$)}
    \begin{enumerate}
        \item Par linéarité de la moyenne :
        \[ \overline{y} = 1,2\overline{x} + 1 = 1,2 \times 9,5 + 1 = 11,4 + 1 = 12,4 \]
        Pour l'écart-type, l'indicateur de dispersion est insensible à la translation ($b=1$) et multiplié par le coefficient d'échelle $a = 1,2$ :
        \[ \sigma(Y) = 1,2 \times \sigma(X) = 1,2 \times 2,4 = 2,88 \]
        \textbf{Interprétation :} La moyenne augmente de près de $3$ points, ce qui rehausse le niveau global. Cependant, l'écart-type augmente également ($\sigma(Y) > \sigma(X)$), ce qui montre que la transformation a « étiré » les notes et accru les écarts mutuels entre les candidats.
        
        \item Le coefficient multiplicateur $a = 1,2$ étant strictement positif, l'ordre des données est conservé. La médiane subit directement la transformation affine :
        \[ M_e(Y) = 1,2 \times M_e(X) + 1 = 1,2 \times 9 + 1 = 11,8 \]
        \textit{Justification rigoureuse :} Par définition, au moins 50\% des individus ont une note $X \le 9$. Comme l'application $t \mapsto 1,2t+1$ est strictement croissante, l'inégalité $X \le 9$ équivaut à $1,2X+1 \le 1,2\times9+1$, c'est-à-dire $Y \le 11,8$. Ainsi, au moins 50\% des individus vérifient $Y \le 11,8$. Le même raisonnement s'applique pour la condition $Y \ge 11,8$.
        
        Pour l'écart interquartile $I_Q(Y)$, les indicateurs de dispersion ne dépendent pas de la translation :
        \[ I_Q(Y) = Q_3(Y) - Q_1(Y) = (1,2Q_3(X)+1) - (1,2Q_1(X)+1) = 1,2 \times I_Q(X) \]
        Or $I_Q(X) = Q_3(X) - Q_1(X) = 11,5 - 7,5 = 4$. Donc :
        \[ I_Q(Y) = 1,2 \times 4 = 4,8 \]
    \end{enumerate}

    \item \textbf{Cas d'un coefficient multiplicateur négatif ($Z = -0,5X + 20$)}
    \begin{enumerate}
        \item \textbf{Démonstration générale :} Soit $Y = aX+b$ avec $a < 0$. On sait que $\overline{y} = a\overline{x}+b$.
        Par définition de la variance :
        \[ V(Y) = \frac{1}{N}\sum_{i=1}^{k} n_i (y_i - \overline{y})^2 \]
        En substituant $y_i = ax_i + b$ et $\overline{y} = a\overline{x} + b$ :
        \[ y_i - \overline{y} = (ax_i + b) - (a\overline{x} + b) = a(x_i - \overline{x}) \]
        En élevant au carré : $(y_i - \overline{y})^2 = a^2 (x_i - \overline{x})^2$. Par linéarité de la somme :
        \[ V(Y) = \frac{1}{N}\sum_{i=1}^{k} n_i a^2 (x_i - \overline{x})^2 = a^2 \left( \frac{1}{N}\sum_{i=1}^{k} n_i (x_i - \overline{x})^2 \right) = a^2 V(X) \]
        En passant à la racine carrée, puisque $\sqrt{a^2} = |a|$ :
        \[ \sigma(Y) = \sqrt{V(Y)} = \sqrt{a^2 V(X)} = |a|\sqrt{V(X)} = |a|\sigma(X) \]
        
        \item Pour la série $Z$, $a = -0,5$. On a donc :
        \[ \sigma(Z) = |-0,5| \times \sigma(X) = 0,5 \times 2,4 = 1,2 \]
        
        \item L'application $t \mapsto -0,5t + 20$ est strictement décroissante car son coefficient directeur est négatif. Par conséquent, elle renverse l'ordre des données. Les plus petites notes initiales deviennent les plus grandes notes transformées et vice-versa. Les rôles des quartiles s'échangent donc :
        \begin{itemize}
            \item $Q_1(Z) = -0,5 \times Q_3(X) + 20 = -0,5 \times 11,5 + 20 = -5,75 + 20 = 14,25$
            \item $Q_3(Z) = -0,5 \times Q_1(X) + 20 = -0,5 \times 7,5 + 20 = -3,75 + 20 = 16,25$
        \end{itemize}
    \end{enumerate}
\end{enumerate}

\vspace{0.5cm}

\section*{Exercice 3 : Fusion de populations et formule de König-Huygens}

\begin{enumerate}
    \item \textbf{Analyse globale de position}
    \begin{enumerate}
        \item L'effectif total de la population fusionnée est $N = N_1 + N_2 = 60 + 40 = 100$. La moyenne globale $\overline{x}$ est donnée par la moyenne pondérée des moyennes de chaque groupe :
        \[ \overline{x} = \frac{N_1\overline{x}_1 + N_2\overline{x}_2}{N_1 + N_2} = \frac{60 \times 25 + 40 \times 35}{100} = \frac{1500 + 1400}{100} = \frac{2900}{100} = 29 \text{ min} \]
        
        \item La moyenne arithmétique simple de $\overline{x}_1$ et $\overline{x}_2$ vaut $\frac{25+35}{2} = 30$ min. La moyenne globale ($29$ min) n'est pas égale à cette moyenne simple car les deux sous-populations possèdent des effectifs différents ($N_1 \neq N_2$). La formule de la moyenne globale montre qu'il s'agit d'un barycentre où chaque sous-population est pondérée par son effectif (le groupe A, plus nombreux, « attire » légitimement la moyenne globale vers sa propre moyenne de $25$ min).
    \end{enumerate}

    \item \textbf{Analyse globale de dispersion}
    \begin{enumerate}
        \item D'après la formule de König-Huygens appliquée au Groupe A, on a $V_1 = M_A - \overline{x}_1^2$, soit $\sigma_1^2 = M_A - \overline{x}_1^2$. On en isole la moyenne des carrés $M_A$ :
        \[ M_A = \sigma_1^2 + \overline{x}_1^2 = 4^2 + 25^2 = 16 + 625 = 641 \]
        
        \item De la même manière pour le Groupe B :
        \[ M_B = \sigma_2^2 + \overline{x}_2^2 = 6^2 + 35^2 = 36 + 1225 = 1261 \]
        
        \item La moyenne des carrés pour la population totale, notée $M_{\text{tot}}$, s'obtient en effectuant la moyenne pondérée des moyennes des carrés de chaque sous-population (car la somme totale des carrés est la somme des carrés du Groupe A et du Groupe B) :
        \[ M_{\text{tot}} = \frac{N_1 M_A + N_2 M_B}{N_1 + N_2} = \frac{60 \times 641 + 40 \times 1261}{100} \]
        \[ M_{\text{tot}} = \frac{38460 + 50440}{100} = \frac{88900}{100} = 889 \]
        
        \item Appliquons de nouveau la formule de König-Huygens sur la population globale pour déterminer la variance globale $V$ :
        \[ V = M_{\text{tot}} - \overline{x}^2 = 889 - 29^2 \]
        Or $29^2 = 841$. Donc :
        \[ V = 889 - 841 = 48 \]
        L'écart-type global $\sigma$ est la racine carrée de la variance globale :
        \[ \sigma = \sqrt{V} = \sqrt{48} = 4\sqrt{3} \approx 6,93 \text{ min} \]
        \textbf{Commentaire :} On remarque que l'écart-type global ($\sigma \approx 6,93$ min) est strictement supérieur à $\sigma_1$ ($4$ min) et à $\sigma_2$ ($6$ min). Cela s'explique par le fait que la dispersion globale prend en compte non seulement la dispersion interne à chaque groupe (variabilité intra-classe), mais s'y ajoute l'écart entre les moyennes des deux groupes (variabilité inter-classe, due au décalage important entre $25$ min et $35$ min).
    \end{enumerate}
\end{enumerate}

\end{document}